\documentclass[12pt]{article}
\usepackage[utf8]{inputenc}
\usepackage[T1]{fontenc}
\usepackage{hyperref}
\title{Wasu Language Resources for K-12 Curriculum Development}
\author{Compiled by Claude Sonnet 5.5 (claude-sonnet-5-5) for Andr\'as Kornai}
\date{}
\begin{document}
\maketitle

\section{Introduction}
Wasu (Wa\c{c}u, Wass\'u; ISO 639-3 wsu; Glottocode wasu1238) is an extinct and unattested, presumed language of the Serra do Azul in Joaquim Gomes, Alagoas, Brazil. It is unclassified. Joshua Project reports that the group's original language has largely been replaced by Portuguese. Fabre's 2005 dictionary says its linguistic affiliation is unknown for lack of data, and raises the possibility that the group was formed from diverse ethnic groups. Do not confuse it with Wasusu, a variety of Southern Nambikwara in Mato Grosso, which is a different language.

\section{Primary Educational Resources}
None exist. No data on the language itself was found.

\section{Linguistic and Reference Materials}
\begin{itemize}
\item Wikipedia on the Wasu language: \url{https://en.wikipedia.org/wiki/Wasu_language}
\item Fabre's 2005 dictionary entry ``Wassu'' (Spanish-language PDF): \url{https://www.ling.fi/Entradas%20diccionario/Dic=Wassu.pdf}
\end{itemize}

\section{Religious and Community Resources}
\begin{itemize}
\item Joshua Project profile of the Wassu in Brazil: \url{https://joshuaproject.net/people_groups/15885}
\end{itemize}

\section{Digital and Audio Resources}
None found.

\section{Curriculum Development Context}
Since the language is unattested, no language-medium curriculum is possible. Any educational work would concern the Wassu community's history and culture, taught in Portuguese, and would need to be developed with the community.

\end{document}